\documentclass[10pt,a4paper]{amsart}
\usepackage[margin=19mm]{geometry}
\usepackage{amsmath,amssymb,mathtools}
\usepackage[scr=rsfs,cal=euler]{mathalfa}
\usepackage{xfrac}
\usepackage[unicode,hidelinks,bookmarks=false]{hyperref}
\newcommand{\ie}{i.e.\ }
\newcommand{\cf}{cf.\ }
\newcommand{\resp}{respectively\ }
\newcommand{\Subsection}{\subsection}
\providecommand{\nobreakdash}{\nobreak\hbox{-}}
\setlength{\emergencystretch}{2em}
\multlinegap0pt
\usepackage{tikz}
\usepackage{xcolor}
\usepackage{amssymb}
\newtheorem{lemma}{Lemma}
\newcommand{\De}{\Delta}
\newcommand{\de}{\delta}
\newcommand{\e}{\epsilon}
\title{Near optimal three-fold additive energy bound for points on convex curves}
\author{Adam Cushman, Ciprian Demeter, Shukun Wu}
\date{Source: arXiv:2608.12316v3 (CC BY 4.0)}
\begin{document}
\maketitle

\noindent\emph{Source and licence.} Near optimal three-fold additive energy bound for points on convex curves. Version arXiv:2608.12316v3 (CC BY 4.0), distributed by arXiv under the Creative Commons Attribution 4.0 International licence, http://creativecommons.org/licenses/by/4.0/, as recorded in the arXiv licence metadata for this exact version and shown on the abstract page https://arxiv.org/abs/2608.12316v3. Full licence notice: \texttt{licenses/arxiv-2608.12316-CC-BY-4.0.txt}.\par\smallskip

\noindent\textit{Editorial scope.} Counted unit: the article's lemma lem:projection-estimate (source0.tex lines 1493--1502). The article's complete original proof of it is delivered with it (source0.tex lines 1504--1586, one \texttt{proof} environment of 83 lines). No mathematics is written, repaired or re-derived: the statement and the proof are the article's own lines, in the article's own notation, and no retained line is substituted or re-flowed.  Retained uncounted as the source-local context the proof needs: lines 1475--1489.  Nothing was omitted from any retained span, so the package carries no omission record.  No retained line contains a citation, so the wrapper carries no bibliography and no bibliography entry.  No retained line contains a figure environment, an image-inclusion command or drawing code, so no image is delivered and the package declares no asset dependency.  Editorial formatting only, in addition: an AMS \texttt{amsart} preamble that restates the article's own declarations and macros used by the retained lines, namely the environments \texttt{lemma} on separate counters, together with the standard packages those lines need.  The retained environments print in this document as Lemma 1 in order of appearance, whereas the article prints the same items with the numbering of its own full document.  The complete native source of the article as distributed in the arXiv e-print for this exact version is bundled unchanged as \texttt{sources/207496/source0.tex}.
\begin{lemma}
    \label{lem:rudnev-stevens-input}
    Let \(A \subset \mathbb{R} \) be a sufficiently large finite set. 
    For any \(X \subset A\), define
    \[
    \sigma_X = |X| / (8 |X+X| \log|A|) , \qquad P(X) = \left\{ y  : r_{X+X} (y) \geq \sigma _X \cdot |X| \right\}, 
    \]
    \[
    R(X) = \{x \in X : \left\lvert (x + X)\cap P(X) \right\rvert \geq 3 |X| /4  \}
    .\]
    Then there exists \(B \subset A \) of size \(\left\lvert B \right\rvert \geq \left\lvert A \right\rvert  / 2\) for which
    \[
    E_2(B) \leq E_2(R(B)) \log \left\lvert A \right\rvert 
    .\]
\end{lemma}

\begin{lemma}\label{lem:projection-estimate}
    Let \(A \subset \mathbb{R} \) be a sufficiently large finite set. Then
    \[
    \left\lvert A \right\rvert ^{12} \ll \left\lvert A-A \right\rvert ^{3}  \cdot E_3(A) \cdot  J_3(A)
    .\]
    Moreover, there exists \(B \subset A\) with \(\left\lvert B \right\rvert \geq \left\lvert A \right\rvert / 2\) such that for any \(\e > 0\)
    \[
    E_2(B) ^{3} \left\lvert A \right\rvert ^{6 - \e} \ll _{\e} \left\lvert A+A \right\rvert ^{2} \cdot E_3(A) ^{2} \cdot J_3(A)
    .\]
\end{lemma}

\begin{proof}
    We first introduce the maps
    \begin{equation}
    \label{eq:projection-maps}
        \pi_{-} (r,a,a') = (r-a, r-a') , \qquad \pi_{+} (r,r',b) = (r+b , r'+b).
    \end{equation}
    If \(Z \subset A ^{3} \), then Cauchy-Schwarz gives
    \begin{equation} 
    \label{eq:cauchy-schwarz-projection}
        \left\lvert Z \right\rvert ^{2} \leq E_3(A) \left\lvert \pi_{\pm } (Z) \right\rvert .
    \end{equation}
    Indeed, by Cauchy-Schwarz
    \[
    \left\lvert Z \right\rvert \leq  \left\lvert \pi_{\pm } (Z) \right\rvert ^{\frac{1}{2} } \# \left\{ (z_1,z_2) \in Z ^{2} : \pi_{\pm } (z_1) = \pi _{\pm } (z_2) \right\} ^{\frac{1}{2} }
    ,\]
    and for either \(\pi _{\pm } \) the rightmost term is bounded by
    \[
    E_3(A) = \# \left\{ (a_{i} )\in A^{6}: a_1 - a_2 = a_3 - a_4 = a_5 - a_6 \right\} 
    .\]
    The proof of the lemma follows by two applications of this inequality.
    
    \bigskip  
    
    For the difference set estimate, put \(\delta = \left\lvert A \right\rvert / (11 \left\lvert A-A \right\rvert )\). We define
    \[
    P = \left\{ x : r_{A-A} (x) \geq \delta \left\lvert A \right\rvert  \right\} , \qquad R = \left\{ r \in A : \left\lvert \left( r-A \right) \cap P \right\rvert  \geq 2 \left\lvert A \right\rvert / \sqrt{11} \right\} 
    .\]
    By the definition of \(P\), we have \(\sum _{x \in P} r_{A-A} (x) \geq 10 \left\lvert A \right\rvert ^{2} / 11\), and hence partitioning this sum by the first coordinate gives \(\left\lvert R \right\rvert \gg \left\lvert A \right\rvert \). 
    For each \(r \in R\), there are \(\geq 4 \left\lvert A \right\rvert ^{2} / 11\) pairs \((a,a') \in A ^{2}\) with \(r-a,r-a' \in P\). 
    By the definition of \(P\), there are also \(\geq 10 \left\lvert A \right\rvert ^{2} / 11\) pairs \((a,a')\) with \(a'-a \in P\), and hence by inclusion-exclusion and the fact that \(|R| \gg |A|\),
    \[
    Z_{-}  = \left\{ (r,a,a') \in R \times  A^{2} : r-a,r-a',a'-a \in P \right\} 
    \]
    satisfies \(\left\lvert Z_{-}  \right\rvert \gg \left\lvert A \right\rvert ^{3} \). 
    Note that 
    \[
    \left\lvert \pi_{-} (Z_{-} ) \right\rvert \leq \# \left\{ p_1 - p_2 = p_3 : p_{i} \in P \right\} 
    ,\]
    and since all \(p \in P\) have \(\geq \delta \left\lvert A \right\rvert \) representations as a difference, we have
    \[
    \left( \delta \left\lvert A \right\rvert  \right) ^{3}  \left\lvert \pi_{-} (Z_{-} ) \right\rvert \leq  \# \{ (a_1 - a_2) - (a_3 - a_4) = (a_5 - a_6) : a_i \in A  \}= J_3(A)
    .\]
    Applying (\ref{eq:cauchy-schwarz-projection}) and using \(|Z_-| \gg |A|^3\) gives
    \[
    \left\lvert A \right\rvert ^{6} \ll \left\lvert Z_{-}  \right\rvert ^{2} \leq E_3(A) \left\lvert \pi_{-} (Z_{-} ) \right\rvert \ll \left( \delta \left\lvert A \right\rvert \right) ^{-3 } E_3(A) J_3(A).
    \]
    Recalling \(\de = |A| / (11 |A-A| )\), we prove the first estimate.
    
    \bigskip
    For the sumset, we begin by choosing \(B \subset A\) as in Lemma \ref{lem:rudnev-stevens-input}. Dyadic pigeonholing for \(E_2(R(B))\) gives \(\Delta \geq 1\) such that, letting 
    \[
    P_{\Delta} = \left\{ x : r_{R(B) - R(B)} (x) \in [\Delta, 2 \Delta) \right\},
    \]
    we have
    \begin{equation} 
    \label{eq:choice-of-dyadic}
        E_2(B) \leq \log \left\lvert A \right\rvert E_2(R(B)) \ll \left( \log \left\lvert A \right\rvert  \right) ^{2} \Delta ^{2} \left\lvert P_{\Delta} \right\rvert  .
    \end{equation}
    Note that there are \(\gg \De |P_{\De}|\) pairs \((r,r') \in R(B) ^2\) with \(r'-r \in P_{\De}\). 
    For each such pair, by the definition of \(R(B)\) and inclusion-exclusion, there are \(\geq |B| /2\) many \(b \in B\) with \(r + b, r' + b \in P(B)\). 
    Therefore, \(\left\lvert Z_{+}  \right\rvert \gg \Delta \left\lvert P_{\Delta}  \right\rvert  \left\lvert B \right\rvert \), where
    \[
        Z_{+} = \left\{ (r,r',b)\in R(B)^{2}\times B : r+b,r'+b \in P(B), r'-r \in P_{\Delta}  \right\} 
    .\]
    If \((p_1,p_2) \in \pi_{+} (Z_{+} )\), then \(p_1,p_2 \in P(B)\), \(p_2-p_1 \in P_{\Delta} \). As before, expanding all of these into sums and differences in \(A\) we obtain 
    \[
        \left( \sigma_{B} \left\lvert B \right\rvert  \right) ^{2} \Delta \left\lvert \pi_{+} (Z_{+} )  \right\rvert \ll \# \left\{ (a_1 + a_2) - (a_3 + a_4) = (a_5 - a_6) : a_{i}\in A \right\} = J_3(A)
    .\]
    Applying (\ref{eq:cauchy-schwarz-projection}) and using \(|Z_+| \gg \De |P_{\De}|\, |B|\) gives
    \begin{equation*}
        \Delta ^{2} \left\lvert P_{\Delta} \right\rvert ^{2} \left\lvert B \right\rvert ^{2} \ll \left\lvert Z_{+}  \right\rvert ^{2} \leq E_3(A) \left\lvert \pi_{+} (Z_{+} ) \right\rvert \ll \Delta ^{-1} (\sigma _B |B| ) ^{-2} E_3(A) J_3(A).
    \end{equation*}
    
    We interpolate this with  $\Delta ^{3} \left\lvert P_{\Delta}  \right\rvert \leq E_3(A)$, obtaining 
    \[
    \De ^6 |P_{\De}|^3 |B| ^2 \ll  (\sigma _B |B|)^{-2} E_3(A)^2 J_3(A)
    .\]
    Finally, use (\ref{eq:choice-of-dyadic}) and recall that \(\sigma_B = |B| / (8 |B+B| \log |A|)\) to obtain
    \[
    E_2(B) ^{3} \left\lvert B \right\rvert ^{6} \ll \left\lvert B+B \right\rvert ^{2} \left( \log \left\lvert A \right\rvert  \right) ^{8} E_3(A)^{2} J_3(A)
    .\]
    The desired result follows upon using \(\left\lvert B+B \right\rvert \leq \left\lvert A+A \right\rvert \), \(\left\lvert B \right\rvert \gg \left\lvert A \right\rvert \), and suppressing powers of \(\log \left\lvert A \right\rvert \) by \(\left\lvert A \right\rvert ^{\e}\).
\end{proof}

\end{document}
